% !TEX spellcheck = en_US



% ~~~~~~~~~~~~~~~~~~~~~~~~~~~~~~~~~~~~~~~ V
% (NC) 2026 Haluk Bingol
% github.com/halukbingol/LaTeX-Templates
%
% Licensees may copy, distribute, display, and perform the work and make derivative works 
% and remixes based on it only for non-commercial purposes. 
% ~~~~~~~~~~~~~~~~~~~~~~~~~~~~~~~~~~~~~~~ A




% =====================================================================
%  Linked Lists -- Junior Level
%  ---------------------------------------------------------------------
%  Built on hblib/Template/hblibTemplate.tex with the libraries in ../../hblib/.
%  Programs and their real outputs are in codeoutput/:
%     codeoutput/<Name>.java   a program (classes without main() are only compiled)
%     codeoutput/<Name>.txt    what it printed, regenerated by:  bash run_examples.sh
% =====================================================================

\documentclass[aspectratio=169,10pt,compress]{beamer}

% ~~~~~~~~~~~~~~~~~~~~~~~~~~~~~~~~~~~~~~~~~~~~~~~~~~~~~~~~~~~~~~~~~ V hblib
	\usepackage{../../hblib/hblibPresentation} % lib presentation: theme, i/N, \hSection, algorithm2e, ...
	\usepackage{../../hblib/hblibCodeoutput} % lib code + output: \lstcode, \lstterminal, ...
	\usepackage{../../hblib/hblibDatastructures} % lib data structures: \hString, \hStack, \hQueue
% ~~~~~~~~~~~~~~~~~~~~~~~~~~~~~~~~~~~~~~~~~~~~~~~~~~~~~~~~~~~~~~~~~ A hblib


\usepackage[T1]{fontenc}
\usepackage{lmodern}
\usepackage{booktabs}

\newcommand{\win}{\textbf{Windows}}
\newcommand{\mac}{\textbf{macOS}}

% long Java lines: wrap instead of running off the frame
\lstset{columns=fixed, breaklines=true, breakatwhitespace=false}

\SetAlCapFnt{\small}
\SetAlCapNameFnt{\small}
\SetKwProg{Fn}{function}{}{end}
\SetKw{KwError}{error}

% ---------------------------------------------------------------------
\title{Linked Lists}
\subtitle{Junior Level: Choosing Between Array and Chain}
\author[Bingol]{Ruhan Alpaydin}

\institute[FCIS]{
    Faculty of Computer and Information Sciences\\
    Yeditepe University
}
\date{
	{\tiny \hbTimeStamp}
}

\begin{document}




% ~~~~~~~~~~~~~~~~~~~~~~~~~~~~~~~~~~~~~~ V frm
\begin{frame}[t,fragile]
  \titlepage
\end{frame}
% ~~~~~~~~~~~~~~~~~~~~~~~~~~~~~~~~~~~~~~ A frm




% ~~~~~~~~~~~~~~~~~~~~~~~~~~~~~~~~~~~~~~ V frm
\begin{frame}[t,fragile] \frametitle{Outline}
  \begin{columns}[T]
    \begin{column}{0.48\textwidth}
      \tableofcontents[sections={1-3}, hideallsubsections]
    \end{column}
    \begin{column}{0.48\textwidth}
      \tableofcontents[sections={4-5}, hideallsubsections]
    \end{column}
  \end{columns}
\end{frame}
% ~~~~~~~~~~~~~~~~~~~~~~~~~~~~~~~~~~~~~~ A frm

% ~~~~~~~~~~~~~~~~~~~~~~~~~~~~~~~~~~~~~~ V
\hSection[Trade]{Array vs Linked}
% ~~~~~~~~~~~~~~~~~~~~~~~~~~~~~~~~~~~~~~ A




% ~~~~~~~~~~~~~~~~~~~~~~~~~~~~~~~~~~~~~~ V frm
\begin{frame}[t,fragile] \frametitle{When each wins}
  \begin{columns}[T]
    \begin{column}{0.48\textwidth}
      \begin{block}{Array wins}
        \small
        \begin{itemize}
          \item access by \textbf{index}: \hCode{a[i]} in $O(1)$
          \item sequential iteration: cache friendly
          \item many small elements: one block, no per-element object
          \item the size is known up front
        \end{itemize}
      \end{block}
    \end{column}
    \begin{column}{0.48\textwidth}
      \begin{block}{Linked list wins}
        \small
        \begin{itemize}
          \item insert or delete at the \textbf{ends}: $O(1)$
          \item insert or delete in the middle, given the node: $O(1)$
          \item no resize copies, ever
          \item the size is unknown at the start
        \end{itemize}
      \end{block}
    \end{column}
  \end{columns}
  \vspace{0.5em}
  \small
  \hRemark{The two rules of thumb}: by \emph{index} $\to$ array; by
  \emph{node} or by \emph{end} $\to$ linked list. Every other choice is
  a measurement.
\end{frame}
% ~~~~~~~~~~~~~~~~~~~~~~~~~~~~~~~~~~~~~~ A frm




% ~~~~~~~~~~~~~~~~~~~~~~~~~~~~~~~~~~~~~~ V frm
\begin{frame}[t,fragile] \frametitle{Benchmark: \texttt{ArrayList} vs \texttt{LinkedList}}
  \begin{columns}[T]
    \begin{column}{0.56\textwidth}
      \lstcode[basicstyle=\ttfamily\tiny]{codeoutput/BenchArrayVsLinked.java}{Java}
    \end{column}
    \begin{column}{0.40\textwidth}
      \lstterminal[basicstyle=\ttfamily\scriptsize]{codeoutput/BenchArrayVsLinked.txt}
      \small
      Both fill the list, then sum the elements \emph{by index}:
      \hCode{list.get(i)} for $i = 0, 1, \dots, n-1$.
    \end{column}
  \end{columns}
  \vspace{0.3em}
\end{frame}
% ~~~~~~~~~~~~~~~~~~~~~~~~~~~~~~~~~~~~~~ A frm




% ~~~~~~~~~~~~~~~~~~~~~~~~~~~~~~~~~~~~~~ V frm
\begin{frame}[t,fragile] \frametitle{Why the second number is so large}
  \begin{center}\small
  \begin{tabular}{@{}lcl@{}}
    \toprule
    \textbf{Container} & \textbf{Call} & \textbf{Cost per call} \\
    \midrule
    \hCode{ArrayList<Integer>}  & \hCode{a.get(i)} & $O(1)$: jump to \hCode{elementData[i]} \\
    \hCode{LinkedList<Integer>} & \hCode{l.get(i)} & $O(i)$: walk \hCode{i} links from \hCode{first} \\
    \bottomrule
  \end{tabular}
  \end{center}
  \vspace{0.4em}
  \small
  So the whole loop is:
  \[
    \text{ArrayList: } O(n); \qquad \text{LinkedList: } O(n^2).
  \]
  \vspace{0.2em}
  \hRemark{The benchmark is not unfair to \texttt{LinkedList}}: it is
  showing the one operation \texttt{LinkedList} is bad at, on purpose.
  The \hCode{get(i)} loop is the wrong workload for a linked list;
  the correct workload is an \hCode{Iterator}, which is $O(n)$ on both.
\end{frame}
% ~~~~~~~~~~~~~~~~~~~~~~~~~~~~~~~~~~~~~~ A frm

% ~~~~~~~~~~~~~~~~~~~~~~~~~~~~~~~~~~~~~~ V
\hSection[Algos]{The Classic Linked-List Algorithms}
% ~~~~~~~~~~~~~~~~~~~~~~~~~~~~~~~~~~~~~~ A




% ~~~~~~~~~~~~~~~~~~~~~~~~~~~~~~~~~~~~~~ V frm
\begin{frame}[t,fragile] \frametitle{Reverse in place}
  \begin{columns}[T]
    \begin{column}{0.56\textwidth}
      \lstcode[basicstyle=\ttfamily\tiny]{codeoutput/ReverseInPlace.java}{Java}
    \end{column}
    \begin{column}{0.40\textwidth}
      \lstterminal[basicstyle=\ttfamily\scriptsize]{codeoutput/ReverseInPlace.txt}
      \small
      Three pointers walk the chain, redirecting each \hCode{next}
      to point at the previous node. No new nodes, no extra list.
      $O(n)$ time, $O(1)$ space.
    \end{column}
  \end{columns}
\end{frame}
% ~~~~~~~~~~~~~~~~~~~~~~~~~~~~~~~~~~~~~~ A frm




% ~~~~~~~~~~~~~~~~~~~~~~~~~~~~~~~~~~~~~~ V frm
\begin{frame}[t,fragile] \frametitle{Find the middle (slow and fast)}
  \begin{columns}[T]
    \begin{column}{0.56\textwidth}
      \lstcode[basicstyle=\ttfamily\tiny]{codeoutput/MiddleElement.java}{Java}
    \end{column}
    \begin{column}{0.40\textwidth}
      \lstterminal[basicstyle=\ttfamily\scriptsize]{codeoutput/MiddleElement.txt}
      \small
      \hCode{slow} advances one node per iteration; \hCode{fast}
      advances two. When \hCode{fast} reaches the end, \hCode{slow} is
      in the middle.
    \end{column}
  \end{columns}
\end{frame}
% ~~~~~~~~~~~~~~~~~~~~~~~~~~~~~~~~~~~~~~ A frm




% ~~~~~~~~~~~~~~~~~~~~~~~~~~~~~~~~~~~~~~ V frm
\begin{frame}[t,fragile] \frametitle{Detect a cycle (Floyd's algorithm)}
  \begin{columns}[T]
    \begin{column}{0.56\textwidth}
      \lstcode[basicstyle=\ttfamily\tiny]{codeoutput/CycleDetect.java}{Java}
    \end{column}
    \begin{column}{0.40\textwidth}
      \lstterminal[basicstyle=\ttfamily\scriptsize]{codeoutput/CycleDetect.txt}
      \small
      Same slow/fast idea: if there is a cycle, the two pointers must
      eventually meet. If \hCode{fast} reaches \hCode{null}, there is
      no cycle.
    \end{column}
  \end{columns}
\end{frame}
% ~~~~~~~~~~~~~~~~~~~~~~~~~~~~~~~~~~~~~~ A frm




% ~~~~~~~~~~~~~~~~~~~~~~~~~~~~~~~~~~~~~~ V frm
\begin{frame}[t,fragile] \frametitle{Why slow and fast always meet}
  \begin{columns}[T]
    \begin{column}{0.5\textwidth}
      \begin{block}{Suppose the cycle has length $c$}
        \small
        Once \hCode{slow} enters the cycle, \hCode{fast} is already
        inside. Each iteration the gap shrinks by one. After at most
        $c$ iterations the gap is zero: they meet.
      \end{block}
    \end{column}
    \begin{column}{0.46\textwidth}
      \begin{block}{Consequence}
        \small
        No extra memory. $O(n)$ time. This is what the JDK uses in
        \hCode{LinkedList}'s internal cycle detection, and what every
        garbage collector does to detect unreachable cycles.
      \end{block}
    \end{column}
  \end{columns}
  \vspace{0.4em}
  \small
  \hRemark{Two pointers, three lines of code}: Floyd's algorithm is one
  of the few $O(1)$-space primitives on a linked list.
\end{frame}
% ~~~~~~~~~~~~~~~~~~~~~~~~~~~~~~~~~~~~~~ A frm




% ~~~~~~~~~~~~~~~~~~~~~~~~~~~~~~~~~~~~~~ V frm
\begin{frame}[t,fragile] \frametitle{Merge two sorted lists}
  \begin{columns}[T]
    \begin{column}{0.56\textwidth}
      \lstcode[basicstyle=\ttfamily\tiny]{codeoutput/MergeSorted.java}{Java}
    \end{column}
    \begin{column}{0.40\textwidth}
      \lstterminal[basicstyle=\ttfamily\scriptsize]{codeoutput/MergeSorted.txt}
      \small
      A dummy node simplifies the boundary: no special case for ``the
      result is empty''. Each step attaches the smaller of the two
      heads. $O(n + m)$.
    \end{column}
  \end{columns}
\end{frame}
% ~~~~~~~~~~~~~~~~~~~~~~~~~~~~~~~~~~~~~~ A frm

% ~~~~~~~~~~~~~~~~~~~~~~~~~~~~~~~~~~~~~~ V
\hSection[Choose]{Choosing the Right Container}
% ~~~~~~~~~~~~~~~~~~~~~~~~~~~~~~~~~~~~~~ A




% ~~~~~~~~~~~~~~~~~~~~~~~~~~~~~~~~~~~~~~ V frm
\begin{frame}[t,fragile] \frametitle{Decision tree}
  \begin{columns}[T]
    \begin{column}{0.48\textwidth}
      \begin{block}{I access by index}
        \small
        \hCode{ArrayList} for general use;\\
        \hCode{int[]} if primitives;\\
        never \hCode{LinkedList}.
      \end{block}
      \begin{block}{I only add or remove at the ends}
        \small
        \hCode{ArrayDeque} for both stack and queue;\\
        \hCode{LinkedList} if you need a \hCode{List} interface too.
      \end{block}
    \end{column}
    \begin{column}{0.48\textwidth}
      \begin{block}{I insert or remove in the middle}
        \small
        \hCode{LinkedList} given the node;\\
        otherwise \hCode{ArrayList} with an index walk.
      \end{block}
      \begin{block}{I iterate over every element}
        \small
        Both are $O(n)$; \hCode{ArrayList} wins on constants \\
        because the CPU cache helps.
      \end{block}
    \end{column}
  \end{columns}
  \vspace{0.4em}
  \small
  \hRemark{Reach for the interface} (\hCode{List}, \hCode{Deque}),
  not the class, so the choice is easy to change.
\end{frame}
% ~~~~~~~~~~~~~~~~~~~~~~~~~~~~~~~~~~~~~~ A frm

% ~~~~~~~~~~~~~~~~~~~~~~~~~~~~~~~~~~~~~~ V
\hSection[JDK]{The JDK's List}
% ~~~~~~~~~~~~~~~~~~~~~~~~~~~~~~~~~~~~~~ A




% ~~~~~~~~~~~~~~~~~~~~~~~~~~~~~~~~~~~~~~ V frm
\begin{frame}[t,fragile] \frametitle{The JDK, in one table}
  \begin{center}\footnotesize
  \begin{tabular}{@{}lccL{0.34\textwidth}@{}}
    \toprule
    \textbf{Operation} & \hCode{ArrayList} & \hCode{LinkedList} & \textbf{Note} \\
    \midrule
    \hCode{get(i)}        & $O(1)$ & $O(i)$          & index vs walk \\
    \hCode{add(x)}        & $O(1)$ amortized & $O(1)$ & append at tail \\
    \hCode{add(0, x)}     & $O(n)$ & $O(1)$          & shift vs relink \\
    \hCode{remove(0)}     & $O(n)$ & $O(1)$          & shift vs relink \\
    \hCode{remove(x)}     & $O(n)$ & $O(n)$          & search first \\
    iterate               & $O(n)$ & $O(n)$          & cache vs pointers \\
    memory per element    & 4 bytes + ref & 24 bytes & header + 3 refs \\
    \bottomrule
  \end{tabular}
  \end{center}
  \vspace{0.3em}
  \small
  \hRemark{The table is the answer}: \hCode{ArrayList} for reads and
  indexed access; \hCode{LinkedList} for queue-like workloads. The JDK
  ships both so you can pick.
\end{frame}
% ~~~~~~~~~~~~~~~~~~~~~~~~~~~~~~~~~~~~~~ A frm

% ~~~~~~~~~~~~~~~~~~~~~~~~~~~~~~~~~~~~~~ V
\hSection[Wrap]{Wrap-up}
% ~~~~~~~~~~~~~~~~~~~~~~~~~~~~~~~~~~~~~~ A




% ~~~~~~~~~~~~~~~~~~~~~~~~~~~~~~~~~~~~~~ V frm
\begin{frame}[t,fragile] \frametitle{Big ideas}
  \begin{itemize}
    \item An \textbf{array} gives $O(1)$ access by index and pays
          $O(n)$ for insert or delete in the middle.
    \item A \textbf{linked list} gives $O(1)$ insert or delete at
          the ends and pays $O(n)$ for access by index.
    \item \hCode{LinkedList.get(i)} is $O(i)$: a loop over it is
          $O(n^2)$, which is why the benchmark on slide~4 is unfair
          to it \emph{on purpose}.
    \item \hRemark{Iterate, do not index}: use the for-each loop or an
          \hCode{Iterator}; both containers are $O(n)$ that way.
    \item Floyd's cycle detection and the slow/fast middle-finding are
          the two $O(1)$-space tricks every linked-list programmer
          should know.
    \item \hCode{ArrayDeque} beats both for pure stack or queue use:
          it is a resizable array with two ends.
  \end{itemize}
\end{frame}
% ~~~~~~~~~~~~~~~~~~~~~~~~~~~~~~~~~~~~~~ A frm




% ~~~~~~~~~~~~~~~~~~~~~~~~~~~~~~~~~~~~~~ V frm
\begin{frame}[t,fragile] \frametitle{Check yourself}
  \begin{enumerate}\small
    \item Why is the \hCode{ArrayList.get(i)} loop $O(n)$ and the
          \hCode{LinkedList.get(i)} loop $O(n^2)$?
    \item When does \hCode{LinkedList.add(0, x)} beat
          \hCode{ArrayList.add(0, x)}?
    \item Why does the benchmark \emph{not} show
          \hCode{LinkedList.removeFirst()} being slow?
    \item What does the \hCode{dummy} node in \hCode{MergeSorted}
          save you?
    \item Why can Floyd's algorithm find a cycle in $O(1)$ memory
          but not the length of the cycle in $O(1)$ memory?
    \item Which class do you reach for when the workload is a pure
          FIFO queue? Which for pure LIFO?
  \end{enumerate}
\end{frame}
% ~~~~~~~~~~~~~~~~~~~~~~~~~~~~~~~~~~~~~~ A frm




% ~~~~~~~~~~~~~~~~~~~~~~~~~~~~~~~~~~~~~~ V frm
\begin{frame}[t,fragile] \frametitle{Practice}
  \begin{enumerate}\small
    \item Re-run the benchmark with the loop replaced by a for-each.
          How do the two containers compare now?
    \item Add \hCode{addLast} to \hCode{ReverseInPlace}'s list, then
          reverse again. Does it still work?
    \item Extend \hCode{CycleDetect} to return the \emph{start} of the
          cycle, not just \hCode{true} or \hCode{false}.
    \item Extend \hCode{MergeSorted} to merge three sorted lists.
    \item Use \hCode{LinkedList} as a stack (push/pop) and as a queue
          (offer/poll). Time each against \hCode{ArrayDeque}.
    \item Write a benchmark that shows \hCode{LinkedList} being
          \emph{faster} than \hCode{ArrayList}. Which operation
          makes it win?
  \end{enumerate}
\end{frame}
% ~~~~~~~~~~~~~~~~~~~~~~~~~~~~~~~~~~~~~~ A frm




% ~~~~~~~~~~~~~~~~~~~~~~~~~~~~~~~~~~~~~~ V frm
\begin{frame}[t,fragile] \frametitle{References}
  \begin{enumerate}\small
    \item Robert Lafore, \emph{Data Structures \& Algorithms in Java},
          2nd ed., Sams, 2002.
    \item Robert Sedgewick and Kevin Wayne, \emph{Algorithms},
          4th ed., Addison-Wesley, 2011. (Chapter 1.3: linked lists
          and array resizing.)
    \item Thomas H. Cormen, Charles E. Leiserson, Ronald L. Rivest,
          Clifford Stein, \emph{Introduction to Algorithms}, 4th ed.,
          MIT Press, 2022. (Cycle detection, reversing, merging.)
    \item Donald E. Knuth, \emph{The Art of Computer Programming},
          Vol.~3: \emph{Sorting and Searching}, 2nd ed.,
          Addison-Wesley, 1998.
    \item Oracle, \emph{Java SE 21 API Specification}:
          \hCode{java.util.ArrayList},
          \hCode{java.util.LinkedList}.
  \end{enumerate}
\end{frame}
% ~~~~~~~~~~~~~~~~~~~~~~~~~~~~~~~~~~~~~~ A frm




% ~~~~~~~~~~~~~~~~~~~~~~~~~~~~~~~~~~~~~~ V frm
\begin{frame}[t,fragile]
  \begin{center}
    \vspace{4cm}
    \hRemark{\Huge Questions?}
  \end{center}
\end{frame}
% ~~~~~~~~~~~~~~~~~~~~~~~~~~~~~~~~~~~~~~ A frm




\end{document}